\section{Особенности реализации}

Программа написана на языке Go. Её основная логика находится в пакете
\texttt{automation}. В программе есть две главные структуры: \texttt{Rule}
хранит правило изменения клеток, а \texttt{Field} --- квадратное поле и
состояние каждой клетки. Состояние хранится как \texttt{bool}: \texttt{true}
означает, что клетка жива, а \texttt{false} --- что мертва. При расчёте следующего шага использует правило из \texttt{Rule}.
\subsection{Структура \texttt{Rule}}

Структура \texttt{Rule} хранит правило, которое программа использует при
обновлении поля. В ней есть один элемент:

\begin{itemize}
    \item \texttt{Number} типа \texttt{uint32} --- число, в котором записано
    правило. Его биты определяют, должна ли клетка стать живой или мёртвой.
\end{itemize}

\subsubsection{Функция \texttt{RuleFromVariant}}

\textbf{Вход:} номер варианта \texttt{variant}.

\textbf{Выход:} объект \texttt{Rule} и пустая ошибка при успехе; нулевое
правило и ошибка, если вариант меньше 1.

\textbf{Описание:} умножает номер варианта на постоянный коэффициент и
формирует правило.

\begin{lstlisting}[style=gostyle,caption={Создание правила по номеру варианта}]
func RuleFromVariant(variant int) (Rule, error) {
    if variant < 1 {
        return Rule{}, fmt.Errorf("variant number must be positive")
    }
    number := uint32(variant) * ruleMultiplier
    if number > math.MaxUint32 {
        return Rule{}, fmt.Errorf("variant number must not exceed %d", math.MaxUint32/ruleMultiplier)
    }
    return Rule{Number: number}, nil
}
\end{lstlisting}

\textbf{Сложность:} $O(1)$ по времени и памяти.

\subsubsection{Метод \texttt{Bits}}
\textbf{Вход:} объект-получатель \texttt{r} типа \texttt{Rule}.

\textbf{Выход:} строка из 32 символов с двоичным номером правила.

\textbf{Описание:} представляет число правила в виде двоичной строки; объект
не изменяется.

\begin{lstlisting}[style=gostyle,caption={Представление правила в двоичном виде}]
func (r Rule) Bits() string {
    return fmt.Sprintf("%032b", r.Number)
}
\end{lstlisting}

\textbf{Сложность:} $O(1)$ по времени и памяти, так как длина строки постоянна.

\subsubsection{Метод \texttt{transition}}
\textbf{Вход:} правило \texttt{r} и номер конфигурации
\texttt{configuration} от 0 до 31.

\textbf{Выход:} \texttt{true}, если центральная клетка должна быть живой,
иначе \texttt{false}.

\textbf{Описание:} проверяет бит правила для заданной конфигурации; конфигурация
0 соответствует старшему биту.

\begin{lstlisting}[style=gostyle,caption={Получение состояния клетки по правилу}]
func (r Rule) transition(configuration int) bool {
    return r.Number&(1<<uint(31-configuration)) != 0
}
\end{lstlisting}

\textbf{Сложность:} $O(1)$ по времени и памяти.

\subsection{Структура \texttt{Field}}

Структура \texttt{Field} хранит квадратное поле и правило для его обновления.
В ней есть два элемента:

\begin{itemize}
    \item \texttt{cells} типа \texttt{[][]bool} --- матрица состояний клеток.
    Значение \texttt{true} означает живую клетку, \texttt{false} --- мёртвую;
    \item \texttt{rule} типа \texttt{Rule} --- правило, используемое для
    вычисления нового состояния поля.
\end{itemize}

Размер поля задаётся количеством строк и столбцов в \texttt{cells}. При
расчёте следующего поколения программа сначала заполняет новую матрицу, а
затем заменяет ею старую.

\subsubsection{Функция \texttt{New}}
\textbf{Вход:} размер поля \texttt{size} и правило \texttt{rule}.

\textbf{Выход:} указатель на созданный объект \texttt{Field}.

\textbf{Описание:} создаёт квадратное поле, изначально заполненное мёртвыми
клетками, и сохраняет в нём правило.

\begin{lstlisting}[style=gostyle,caption={Создание пустого поля}]
func New(size int, rule Rule) *Field {
    return &Field{cells: makeCells(size), rule: rule}
}
\end{lstlisting}

\textbf{Сложность:} $O(n^2)$ по времени и памяти.

\subsubsection{Метод \texttt{Size}}
\textbf{Вход:} объект-получатель \texttt{f} типа \texttt{*Field}.

\textbf{Выход:} целое число --- длина стороны поля.

\textbf{Описание:} возвращает количество строк поля, не изменяя его.

\begin{lstlisting}[style=gostyle,caption={Получение размера поля}]
func (f *Field) Size() int {
    return len(f.cells)
}
\end{lstlisting}

\textbf{Сложность:} $O(1)$ по времени и памяти.

\subsubsection{Метод \texttt{At}}
\textbf{Вход:} объект \texttt{f} типа \texttt{*Field}, строка
\texttt{row} и столбец \texttt{col}.

\textbf{Выход:} состояние клетки типа \texttt{bool}.

\textbf{Описание:} возвращает состояние клетки; координаты за границей поля
переходят на противоположную сторону.

\begin{lstlisting}[style=gostyle,caption={Чтение состояния клетки}]
func (f *Field) At(row, col int) bool {
    return f.cells[f.wrap(row)][f.wrap(col)]
}
\end{lstlisting}

\textbf{Сложность:} $O(1)$ по времени и памяти.

\subsubsection{Функция \texttt{makeCells}}
\textbf{Вход:} размер матрицы \texttt{size}.

\textbf{Выход:} квадратная матрица типа \texttt{[][]bool}, заполненная
значениями \texttt{false}.

\textbf{Описание:} выделяет память для строк матрицы и создаёт пустое поле.

\begin{lstlisting}[style=gostyle,caption={Создание матрицы клеток}]
func makeCells(size int) [][]bool {
    cells := make([][]bool, size)
    for row := range cells {
        cells[row] = make([]bool, size)
    }
    return cells
}
\end{lstlisting}

\textbf{Сложность:} $O(n^2)$ по времени и памяти.

\subsubsection{Метод \texttt{Set}}
\textbf{Вход:} объект поля \texttt{f}, координаты \texttt{row} и
\texttt{col}, новое состояние \texttt{state}.

\textbf{Выход:} отдельно не возвращается; обновляется объект \texttt{Field}.

\textbf{Описание:} записывает заданное состояние в выбранную клетку.

\begin{lstlisting}[style=gostyle,caption={Задание состояния клетки}]
func (f *Field) Set(row, col int, state bool) {
    f.cells[f.wrap(row)][f.wrap(col)] = state
}
\end{lstlisting}

\textbf{Сложность:} $O(1)$ по времени и памяти.

\subsubsection{Метод \texttt{Toggle}}
\textbf{Вход:} объект \texttt{f} типа \texttt{*Field} и координаты
клетки \texttt{row}, \texttt{col}.

\textbf{Выход:} отдельно не возвращается; обновляется объект \texttt{Field}.

\textbf{Описание:} меняет состояние клетки на противоположное.

\begin{lstlisting}[style=gostyle,caption={Переключение состояния клетки}]
func (f *Field) Toggle(row, col int) {
    row, col = f.wrap(row), f.wrap(col)
    f.cells[row][col] = !f.cells[row][col]
}
\end{lstlisting}

\textbf{Сложность:} $O(1)$ по времени и памяти.

\subsubsection{Метод \texttt{Clear}}
\textbf{Вход:} объект-получатель \texttt{f} типа \texttt{*Field};
дополнительных аргументов нет.

\textbf{Выход:} отдельно не возвращается; все клетки объекта \texttt{Field}
становятся мёртвыми.

\textbf{Описание:} проходит по матрице и сбрасывает состояние каждой клетки.

\begin{lstlisting}[style=gostyle,caption={Очистка поля}]
func (f *Field) Clear() {
    for row := range f.cells {
        for col := range f.cells[row] {
            f.cells[row][col] = false
        }
    }
}
\end{lstlisting}

\textbf{Сложность:} $O(n^2)$ по времени и $O(1)$ дополнительной памяти.

\subsubsection{Метод \texttt{Randomize}}
\textbf{Вход:} объект \texttt{f} типа \texttt{*Field} и начальное
значение генератора \texttt{seed}.

\textbf{Выход:} отдельно не возвращается; заполняется объект \texttt{Field}.

\textbf{Описание:} случайно задаёт состояние каждой клетки. При одинаковых
размере поля и \texttt{seed} результат повторяется.

\begin{lstlisting}[style=gostyle,caption={Случайное заполнение поля}]
func (f *Field) Randomize(seed int64) {
    rng := rand.New(rand.NewSource(seed))
    for row := range f.cells {
        for col := range f.cells[row] {
            f.cells[row][col] = rng.Intn(2) == 1
        }
    }
}
\end{lstlisting}

\textbf{Сложность:} $O(n^2)$ по времени и $O(1)$ дополнительной памяти.

\subsubsection{Метод \texttt{neighbourhoodIndex}}
\textbf{Вход:} объект \texttt{f} типа \texttt{*Field} и координаты
центральной клетки \texttt{row}, \texttt{col}.

\textbf{Выход:} номер конфигурации от 0 до 31.

\textbf{Описание:} кодирует состояния пяти клеток в порядке: центр, верх,
право, низ, лево.

\begin{lstlisting}[style=gostyle,caption={Кодирование конфигурации соседей}]
func (f *Field) neighbourhoodIndex(row, col int) int {
    bits := [5]bool{
        f.At(row, col),
        f.At(row-1, col),
        f.At(row, col+1),
        f.At(row+1, col),
        f.At(row, col-1),
    }
    index := 0
    for _, bit := range bits {
        index <<= 1
        if bit {
            index++
        }
    }
    return index
}
\end{lstlisting}

\textbf{Сложность:} $O(1)$ по времени и памяти.

\subsubsection{Метод \texttt{wrap}}
\textbf{Вход:} объект \texttt{f} типа \texttt{*Field} и координата
\texttt{value}.

\textbf{Выход:} индекс от 0 до \texttt{size - 1}, где \texttt{size} --- размер
поля.

\textbf{Описание:} приводит координату к полю с учётом перехода через границы.

\begin{lstlisting}[style=gostyle,caption={Приведение координаты к границам поля}]
func (f *Field) wrap(value int) int {
    size := f.Size()
    return (value%size + size) % size
}
\end{lstlisting}

\textbf{Сложность:} $O(1)$ по времени и памяти.

\subsubsection{Метод \texttt{Step}}
\textbf{Вход:} объект-получатель \texttt{f} типа \texttt{*Field};
дополнительных аргументов нет.

\textbf{Выход:} отдельно не возвращается; текущее состояние поля заменяется
следующим поколением.

\textbf{Описание:} вычисляет состояние всех клеток во временной матрице и
только затем заменяет ею текущее поле. Поэтому при расчёте одного поколения
все клетки используют состояния предыдущего поколения.

\begin{lstlisting}[style=gostyle,caption={Вычисление следующего поколения}]
func (f *Field) Step() {
    next := makeCells(f.Size())
    for row := range f.cells {
        for col := range f.cells[row] {
            next[row][col] = f.rule.transition(f.neighbourhoodIndex(row, col))
        }
    }
    f.cells = next
}
\end{lstlisting}

\textbf{Сложность:} $O(n^2)$ по времени и $O(n^2)$ дополнительной памяти.
